%% ============================================================
%%  Lab 4 Report – Motor Control using RTOS
%%  Embedded Systems Development  04-633 A
%%  Carnegie Mellon University Africa  –  April 2026
%% ============================================================
\documentclass[11pt, a4paper]{article}

% ── Packages ─────────────────────────────────────────────────
\usepackage[top=2.5cm, bottom=2.5cm, left=2.8cm, right=2.8cm]{geometry}
\usepackage{graphicx}
\usepackage{xcolor}
\usepackage{hyperref}
\usepackage{booktabs}
\usepackage{array}
\usepackage{enumitem}
\usepackage{titlesec}
\usepackage{parskip}
\usepackage{caption}
\usepackage{fancyhdr}
\usepackage{microtype}
\usepackage{float}
\usepackage{amsmath}
\usepackage{setspace}

% ── Colours ──────────────────────────────────────────────────
\definecolor{cmuRed}{RGB}{196, 18, 48}
\definecolor{cmuGray}{RGB}{90, 90, 90}

% ── Hyperref ─────────────────────────────────────────────────
\hypersetup{colorlinks=true, linkcolor=black, urlcolor=cmuRed, citecolor=black}

% ── Section formatting – clean, no red rule ──────────────────
\titleformat{\section}{\large\bfseries}{}{0em}{}
\titleformat{\subsection}{\normalsize\bfseries\itshape}{}{0em}{}
\titlespacing{\section}{0pt}{10pt}{4pt}
\titlespacing{\subsection}{0pt}{7pt}{2pt}

% ── Header / Footer ──────────────────────────────────────────
\pagestyle{fancy}
\fancyhf{}
\renewcommand{\headrulewidth}{0.4pt}
\lhead{\small\color{cmuGray} 04-633 A \ Embedded Systems Development}
\rhead{\small\color{cmuGray} Lab 4: Motor Control using RTOS}
\cfoot{\small\thepage}

% ── Table column types ───────────────────────────────────────
\newcolumntype{L}[1]{>{\raggedright\arraybackslash}p{#1}}
\newcolumntype{C}[1]{>{\centering\arraybackslash}p{#1}}

%% ============================================================
\begin{document}
\pagenumbering{gobble}

%% ── COVER PAGE ──────────────────────────────────────────────
\begin{titlepage}
  \centering

  \vspace*{0.6cm}
  \includegraphics[width=3.8cm]{logo.png}\\[0.9cm]

  {\fontsize{21}{25}\selectfont\bfseries\color{cmuRed}
    Carnegie Mellon University Africa}\\[0.3cm]
  {\large\itshape Embedded Systems Development}\\[0.1cm]
  {\normalsize 04-633 A}\\[2.0cm]

  {\Large\bfseries Lab 4: Motor Control using RTOS}\\[0.35cm]
  {\normalsize\color{cmuRed} FreeRTOS, Hardware PWM \& Servo Control on STM32}\\[2.4cm]

  \begin{tabular}{rl}
    \textbf{Name:}      & Abba Uba Said \\[0.25em]
    \textbf{Andrew ID:} & \href{mailto:asaiduba@andrew.cmu.edu}{asaiduba@andrew.cmu.edu} \\[0.5em]
    \textbf{Partner:}   & Jeremiah Dosu \\[0.25em]
    \textbf{Andrew ID:} & \href{mailto:jdosu@andrew.cmu.edu}{jdosu@andrew.cmu.edu} \\
  \end{tabular}

  \vfill

  \begin{tabular}{rl}
    \textbf{Course Instructor:}  & Emmanuel Ndashimye \\[0.2em]
    \textbf{Teaching Assistant:} & Millicent Malinga \\[0.2em]
    \textbf{Due Date:}           & April 10, 2026 \\[0.2em]
    \textbf{Demo Date:}          & April 13, 2026 \\
  \end{tabular}

  \vspace{1.4cm}
\end{titlepage}

%% ── Content pages ───────────────────────────────────────────
\pagenumbering{arabic}
\setcounter{page}{1}
\setlength{\parskip}{4pt}

%% ============================================================
\section{Overview of the Implementation}
%% ============================================================

This project implements a real-time servo motor controller on the STM32 Nucleo-F103RB
using FreeRTOS for concurrent task management and hardware PWM via TIM2 for precise
motor control. Users enter a desired servo angle (0--180°) on a 4×4 matrix keypad; live
feedback appears on a 16×2 I\textsuperscript{2}C LCD; and the SG90 servo moves to the
specified position — all running as independent, synchronised RTOS tasks.

\subsection{FreeRTOS Task Design}

\begin{table}[H]
\centering\small
\caption{FreeRTOS Task Summary}
\begin{tabular}{L{2.6cm} C{2.2cm} C{1.8cm} L{6.0cm}}
\toprule
\textbf{Task} & \textbf{Priority} & \textbf{Stack} & \textbf{Responsibility} \\
\midrule
\texttt{KeypadTask}  & Normal       & 256 words & Scans keypad; buffers digits; sends angle on \texttt{\#}; backspace on \texttt{*} \\
\texttt{ServoTask}   & AboveNormal  & 256 words & Reads confirmed angle; updates TIM2 compare register (50\,Hz PWM) \\
\texttt{LCDTask}     & Normal       & 256 words & Row\,0: live typing with blinking cursor; Row\,1: confirmed angle \\
\texttt{defaultTask} & Low          & 128 words & Toggles Nucleo LD2 (PA5) every 3\,s — RTOS heartbeat \\
\midrule
\multicolumn{4}{l}{\small Shared state: \texttt{targetAngle} and \texttt{inputDisplay[5]} — both protected by \texttt{angleMutex}.}\\
\bottomrule
\end{tabular}
\end{table}

\subsection{Hardware Configuration}

The system clock runs at \textbf{72\,MHz} (HSE 8\,MHz $\times$ PLL\,$\times$9).
TIM4 serves as the HAL timebase (separate from the RTOS SysTick), while TIM2 Channel\,1
generates the servo PWM on PA0. I\textsuperscript{2}C1 is remapped to PB8/PB9 for the LCD.

\begin{table}[H]
\centering\small
\caption{Pin Assignment}
\begin{tabular}{C{1.6cm} C{2.8cm} L{8.0cm}}
\toprule
\textbf{Pin(s)} & \textbf{Function} & \textbf{Purpose} \\
\midrule
PA0      & TIM2\_CH1 PWM   & SG90 signal — 50\,Hz, 500–2500\,µs pulse \\
PA5      & GPIO Output     & Nucleo LD2 heartbeat LED (3\,s toggle) \\
PB0--PB3 & GPIO Output     & Keypad rows (driven LOW one at a time to scan) \\
PB4--PB7 & GPIO Input PU   & Keypad columns (read LOW when key pressed) \\
PB8/PB9  & I\textsuperscript{2}C1 SCL/SDA & LCD (HD44780 + PCF8574 backpack, addr 0x27) \\
PD0/PD1  & RCC OSC IN/OUT  & 8\,MHz HSE external crystal \\
\bottomrule
\end{tabular}
\end{table}

\subsection{Servo PWM Calculation}

The SG90 requires a 50\,Hz (20\,ms) PWM signal. With the 72\,MHz APB1 timer clock:
\[
  f_\text{timer} = \frac{72{,}000{,}000}{72} = 1\,\text{MHz}\,(1\,\mu\text{s/tick}),
  \qquad
  T_\text{PWM} = 20{,}000 \times 1\,\mu\text{s} = 20\,\text{ms} = 50\,\text{Hz}
\]
\[
  \text{Pulse}(\theta) = 500 + \frac{\theta \times 2000}{180}\;\text{ticks},
  \quad \theta \in [0°,\;180°]
\]

%% ============================================================
\section{How FreeRTOS and Timer Interrupts Were Used}
%% ============================================================

\subsection{Task Scheduling and Synchronisation}

FreeRTOS is configured with a 1\,kHz tick rate (\texttt{configTICK\_RATE\_HZ = 1000}),
8\,KB heap, and preemptive scheduling. \texttt{ServoTask} holds the highest priority,
guaranteeing prompt PWM updates whenever a new angle is confirmed.
All tasks use \texttt{osDelay()} rather than busy-waiting, yielding the CPU between
iterations: \texttt{KeypadTask} polls every 10\,ms, \texttt{LCDTask} refreshes every 150\,ms,
and \texttt{ServoTask} updates at 20\,ms intervals.

A single mutex (\texttt{angleMutex}) protects both shared variables.
Every task follows the same pattern — acquire mutex, read/write, release immediately —
minimising lock contention and avoiding priority inversion.

\subsection{Hardware Timer Interrupts}

\textbf{TIM4} generates 1\,ms HAL timebase interrupts independently of the RTOS SysTick,
a mandatory separation on STM32 when running FreeRTOS alongside the HAL.

\textbf{TIM2} operates in hardware PWM mode (TIM\_OCMODE\_PWM1) on Channel\,1 (PA0).
Once started with \texttt{HAL\_TIM\_PWM\_Start()}, the timer peripheral autonomously
maintains the 50\,Hz waveform without software intervention.
\texttt{ServoTask} only calls \texttt{\_\_HAL\_TIM\_SET\_COMPARE()} to change the pulse
width, keeping CPU overhead minimal.

%% ============================================================
\section{Challenges, Solutions, and Key Lessons Learned}
%% ============================================================

\subsection{Challenge 1 — Servo Vibrating: Wrong PWM Frequency}

The servo vibrated continuously instead of moving to a set position.
A logic analyser revealed the PWM frequency was \textbf{5\,000\,Hz} — 100× too high.
The CubeMX configuration had \texttt{Prescaler = 719}, \texttt{Period = 19}, producing
a 100\,kHz timer clock and a 0.2\,ms period. The SG90 interprets out-of-range frequencies
as invalid commands, causing uncontrolled coil energisation.

\textit{Solution:} Recalculated from first principles to \texttt{Prescaler = 71},
\texttt{Period = 19999} (1\,MHz tick, 20\,ms period). The servo settled correctly on the
first test after this change.

\subsection{Challenge 2 — Blank LCD: I\textsuperscript{2}C Misconfiguration}

We spent the entire morning in the lab with a completely blank display.
CubeMX had remapped I\textsuperscript{2}C1 to the alternate pins (PB8/PB9)
and enabled \texttt{\_\_HAL\_AFIO\_REMAP\_I2C1\_ENABLE()}, but the physical wires
remained on the default pins (PB6/PB7). Additionally, \texttt{HAL\_MspInit()} called
\texttt{\_\_HAL\_AFIO\_REMAP\_SWJ\_DISABLE()}, which silently killed both JTAG and SWD,
locking out the ST-Link debugger and making every reflash require manual bootloader entry.

\textit{Solution:} Moved SDA/SCL wires to \textbf{PB9/PB8}; changed
\texttt{SWJ\_DISABLE} to \texttt{SWJ\_NOJTAG} to restore ST-Link while freeing PB3/PB4
for keypad use; verified the PCF8574 I\textsuperscript{2}C address (0x27) with a scanner.
The LCD initialised immediately on the next boot.

\subsection{Challenge 3 — Keypad Row 4 Dead: JTAG Pin Conflict}

After fixing the LCD, the entire bottom row (\texttt{*}, \texttt{0}, \texttt{\#}, \texttt{D})
was unresponsive. Row\,4 is driven by PB3 (JTAG TDI pin). The AFIO remap had not taken
effect before GPIO initialisation because \texttt{SWJ\_DISABLE} was overriding it.

\textit{Solution:} \texttt{SWJ\_NOJTAG} — called inside \texttt{HAL\_MspInit()},
which executes before \texttt{MX\_GPIO\_Init()} — releases PB3 and PB4 to the GPIO
peripheral. All 16 keys subsequently responded correctly.

\subsection{Challenge 4 — Empty Task Bodies: No System Output}

The CubeMX-generated skeleton contained only \texttt{osDelay(1)} placeholders.
After flashing, nothing worked: no LED blink, no LCD output, servo stationary.
Additionally, \texttt{servo.c} was missing entirely — declared in \texttt{servo.h}
but never created — causing a linker error on the first build.

\textit{Solution:} All three application tasks were fully implemented, and \texttt{servo.c}
was written from scratch with \texttt{Servo\_Init()} and \texttt{Servo\_SetAngle()}.

\subsection{Results}

After resolving all four challenges, the system operated as expected.
Typing digits on the keypad displayed them in real time on LCD row\,0 with a blinking
cursor; pressing \texttt{\#} committed the angle to \texttt{ServoTask} via the mutex,
moved the servo smoothly, and updated LCD row\,1. The LD2 LED blinked every 3\,s
confirming the RTOS heartbeat.

\begin{figure}[H]
  \centering
  % Automatically loads "result.png" or "result.jpg" if uploaded to Overleaf,
  % otherwise falls back to the placeholder box.
  \IfFileExists{result.png}{
    \includegraphics[width=0.85\textwidth]{result.png}
  }{
    \IfFileExists{result.jpg}{
      \includegraphics[width=0.85\textwidth]{result.jpg}
    }{
      \fbox{\begin{minipage}{0.85\textwidth}\centering\vspace{3.2cm}
        {\color{cmuGray}\itshape
          [Insert photo: Upload ``result.png'' or ``result.jpg'' to Overleaf\\
           It should show the live LCD feedback, servo, and keypad]}
      \vspace{3.2cm}\end{minipage}}
    }
  }
  \caption{Working system: live LCD feedback and SG90 servo responding to keypad input.}
\end{figure}

\subsection{Key Lessons Learned}

\begin{enumerate}[noitemsep, topsep=2pt]
  \item \textbf{Verify timer arithmetic independently.} Never rely on auto-generated
        prescaler values without confirming $f_\text{tick}$ and $T_\text{period}$ on paper.
  \item \textbf{Pin remapping must match physical wiring.} A mismatch between CubeMX
        alternate-function remapping and the actual wire connections produces silent failures.
  \item \textbf{Preserve SWD during development.} \texttt{SWJ\_DISABLE} eliminates the
        only practical reflash path on Nucleo boards and must never be used while debugging.
  \item \textbf{Incremental integration.} Testing each peripheral in isolation before
        combining all tasks would have caught the empty-body and missing-file issues earlier.
  \item \textbf{Mutex discipline.} Acquiring, accessing, and releasing the mutex as quickly
        as possible in every task prevents deadlocks and keeps real-time response predictable.
\end{enumerate}

%% ============================================================
\section{Team Collaboration}
%% ============================================================

This lab was completed in partnership with my teammate, \textbf{Jeremiah}, in accordance
with the course collaboration policy. I focused primarily on the FreeRTOS task architecture,
the servo driver (\texttt{servo.c}), and the TIM2 PWM configuration, while Jeremiah handled
the keypad scanning logic and the LCD driver integration. That said, every major design
decision, pin assignment, and timer calculation was discussed and validated by both of us
before any code was committed to the shared repository.

\subsection{In-Lab Debugging Session}

Jeremiah and I arrived at the lab early in the morning with what I thought was a solid,
nearly-complete implementation. We had discussed the design the night before and I was
confident it would work on the first flash. It did not.

\textbf{The servo was vibrating the moment we powered up.}
Before Jeremiah and I even tried the LCD or keypad, the motor was buzzing continuously
and shaking at every angle I commanded. My first instinct was a wiring problem, so we
spent time tracing every connection on the breadboard --- everything looked correct.
I then connected a logic analyser to PA0 and measured the PWM frequency: it was
\textbf{5\,000\,Hz}, not 50\,Hz --- a factor of 100 off. I traced the error to the timer
values I had written in CubeMX: \texttt{Prescaler = 720 - 1} and \texttt{Period = 20 - 1}
produced a 100\,kHz timer clock with a 0.2\,ms period. I recalculated from the SG90
datasheet on paper, corrected to \texttt{Prescaler = 72 - 1} and \texttt{Period = 20000 - 1},
reflashed the board, and the servo moved cleanly to the target position for the first time.

\textbf{The LCD stayed completely blank for hours.}
With the servo working, Jeremiah and I turned to the display. No matter what we tried,
it showed nothing --- not even garbage characters. Jeremiah verified the
I\textsuperscript{2}C address with a scanner sketch (confirmed 0x27), we inspected the
PCF8574 solder jumpers, and I cross-checked the \texttt{LCD\_Init()} sequence against the
HD44780 datasheet line by line. Still nothing. The breakthrough came when I opened the
STM32CubeIDE pin view and noticed that I\textsuperscript{2}C1 was remapped to PB8 and
PB9 in the configuration --- but our physical wires were connected to PB6 and PB7, the
non-remapped default pins. I moved the SDA wire to PB9 and the SCL wire to PB8 on the
Nucleo header. The LCD displayed ``\texttt{Servo Control}'' on the very next power cycle.
That moment was a genuine relief after hours of confusion.

\textbf{Losing the ST-Link connection mid-session.}
At one point during the day, a reflash attempt failed and the board completely vanished
from STM32CubeIDE. Jeremiah and I tried different cables and ports. I eventually realised
that \texttt{\_\_HAL\_AFIO\_REMAP\_SWJ\_DISABLE()} in \texttt{HAL\_MspInit()} had
silently disabled both JTAG and SWD after the first successful flash. To recover, I held
the BOOT0 pin HIGH while powering on to enter bootloader mode and reflash the firmware.
I immediately changed the call to \texttt{SWJ\_NOJTAG}, which preserves the SWD
interface (used by ST-Link) while freeing the JTAG-only pins (PB3, PB4) for the keypad.

\textbf{Working through it together.}
Throughout the day, Jeremiah kept me grounded when I got tunnel-vision on a single
hypothesis, and I helped him debug the keypad row~4 issue when he ran into the JTAG
pin conflict. We used a shared GitHub Classroom repository with feature branches ---
I maintained \texttt{feature/servo-pwm} and Jeremiah worked on \texttt{feature/lcd-keypad}
--- merging to \texttt{main} only after hardware tests confirmed each feature was stable.
Despite the long and at times frustrating day, working through every bug with Jeremiah
gave me a much deeper understanding of STM32 peripheral configuration and FreeRTOS
integration than any lecture alone could have provided.

%% ============================================================
\section*{Conclusion}
%% ============================================================

This lab provided practical experience with FreeRTOS task scheduling, hardware PWM
generation, I\textsuperscript{2}C communication, and GPIO matrix scanning on the STM32
platform. The debugging process reinforced that systematic hardware verification —
clock trees, pin remapping, and signal measurement — combined with incremental software
integration are essential disciplines for reliable embedded systems development.
The final system meets all lab requirements: concurrent RTOS tasks, hardware PWM servo
control, live LCD feedback, a mutex-protected shared state, and a visible RTOS heartbeat.

%% ============================================================
\section*{References}
%% ============================================================

\begin{enumerate}[noitemsep, topsep=2pt, label={[\arabic*]}]
  \item STMicroelectronics, \textit{STM32F103xB Datasheet}, Rev.\,17, 2022.
        \url{https://www.st.com/resource/en/datasheet/stm32f103rb.pdf}
  \item STMicroelectronics, \textit{STM32CubeIDE User Guide (UM2553)}, 2023.
        \url{https://www.st.com/resource/en/user_manual/um2553-stm32cubeide-user-guide-stmicroelectronics.pdf}
  \item Carnegie Mellon University Africa, \textit{04-633 A Reference Manual}, Canvas, 2026.
  \item STMicroelectronics, \textit{Nucleo-64 User Manual (UM1724)}, Rev.\,14, 2023.
        \url{https://www.st.com/resource/en/user_manual/um1724-stm32-nucleo64-boards-mb1136-stmicroelectronics.pdf}
  \item Hitachi, \textit{HD44780U LCD Controller Datasheet}, 1998.
        \url{https://www.sparkfun.com/datasheets/LCD/HD44780.pdf}
  \item NXP Semiconductors, \textit{PCF8574 Remote 8-Bit I/O Expander Datasheet}, Rev.\,5, 2021.
        \url{https://www.nxp.com/docs/en/data-sheet/PCF8574_PCF8574A.pdf}
  \item Tower Pro, \textit{SG90 Micro Servo Motor Datasheet}.
        \url{http://www.towerpro.com.tw/product/sg90-7/}
\end{enumerate}

\end{document}
